% =====================================================================
%  chapters/chapter-template.tex
%  Kerangka (skeleton) bab modul. Salin file ini menjadi
%  chapters/bab-XX.tex lalu isi konten. Jangan tulis konten penuh di sini.
%
%  Konvensi penamaan bab: chapters/bab-01.tex, bab-02.tex, dst.
%  Daftarkan di main.tex dengan \input{chapters/bab-XX.tex}.
% =====================================================================

\chapter{Pengantar Desain Web}
\label{chap:pengantar}

% ---------------------------------------------------------------------
% 1. CAPAIAN PEMBELAJARAN
% ---------------------------------------------------------------------
\begin{capaian}
Setelah menyelesaikan bab ini, mahasiswa diharapkan mampu:
\begin{itemize}
  \item menjelaskan konsep dasar \emph{World Wide Web} dan cara kerja HTTP;
  \item membedakan peran HTML, CSS, dan JavaScript dalam pengembangan web;
  \item membuat halaman web statis sederhana dengan struktur HTML yang valid;
  \item menerapkan gaya dasar CSS pada elemen halaman.
\end{itemize}
\end{capaian}

% ---------------------------------------------------------------------
% 2. TUJUAN PRAKTIK
% ---------------------------------------------------------------------
\begin{tujuanpraktik}
Pada sesi praktik ini, mahasiswa akan:
\begin{itemize}
  \item menyiapkan lingkungan kerja (editor teks dan peramban);
  \item membuat berkas HTML pertama dan membukanya di peramban;
  \item memvalidasi struktur dokumen HTML.
\end{itemize}
\end{tujuanpraktik}

% ---------------------------------------------------------------------
% 3. ISI MATERI (uraian singkat — tulis konten penuh di sini)
% ---------------------------------------------------------------------
\section{Konsep Dasar}
\label{sec:konsep-dasar}

Tulis paragraf pengantar materi di sini. Gunakan \emph{italic} untuk istilah
asing dan \texttt{monospace} untuk nama berkas atau perintah.

\subsection{Sub-bagian Materi}
\label{subsec:sub-bagian}

Tulis uraian materi di sini.

% ---------------------------------------------------------------------
% 4. LANGKAH KERJA (praktik)
% ---------------------------------------------------------------------
\section{Praktik: Membuat Halaman Pertama}
\label{sec:praktik-pertama}

\begin{langkah}
  \item Buka editor teks dan buat berkas baru bernama \texttt{index.html}.
  \item Ketik kerangka dokumen HTML berikut.
  \item Simpan berkas dan buka di peramban.
  \item Amati hasilnya dan bandingkan dengan struktur kode.
\end{langkah}

% --- Contoh blok kode (listings sudah dimuat ElegantBook) ---
\begin{lstlisting}[language=HTML, caption={Struktur dasar HTML}, label={lst:html-dasar}]
<!DOCTYPE html>
<html lang="id">
<head>
  <meta charset="UTF-8">
  <title>Halaman Pertama</title>
</head>
<body>
  <h1>Halo, Dunia!</h1>
</body>
</html>
\end{lstlisting}

% ---------------------------------------------------------------------
% 5. TANGKAPAN LAYAR (placeholder non-gambar)
%    Input: {jalur aset}{keterangan}{instruksi mahasiswa}
% ---------------------------------------------------------------------
\begin{tangkapanlayar}
  {assets/images/bab-01/hasil-halaman-pertama.png}
  {Hasil halaman pertama yang ditampilkan di peramban.}
  {Buka \texttt{index.html} di peramban, lalu bandingkan tampilan Anda
   dengan contoh pada gambar.}
\end{tangkapanlayar}

% ---------------------------------------------------------------------
% 6. CATATAN / TIPS / PERINGATAN
% ---------------------------------------------------------------------
\begin{catatan}
Tulis catatan penting di sini, misalnya perbedaan antara tag blok dan tag sebaris.
\end{catatan}

\begin{tip}
Tulis kiat praktis di sini, misalnya penggunaan pintasan editor.
\end{tip}

\begin{peringatan}
Tulis peringatan di sini, misalnya kesalahan umum yang sering terjadi.
\end{peringatan}

% ---------------------------------------------------------------------
% 7. LATIHAN
% ---------------------------------------------------------------------
\section{Latihan}
\label{sec:latihan}

\begin{latihan}
  \item Buat halaman HTML berisi judul, paragraf, dan daftar.
  \item Tambahkan tautan ke halaman lain.
  \item Validasi berkas Anda menggunakan validator HTML.
\end{latihan}

% ---------------------------------------------------------------------
% 8. RANGKUMAN (opsional)
% ---------------------------------------------------------------------
\section{Rangkuman}
\label{sec:rangkuman}

Tulis ringkasan poin-poin penting bab ini.